\section{Conclusion}
\label{sec:conclusion}

We investigated whether alignment diminishes a language model's stylistic
capabilities and whether logit-space distribution arithmetic can recover them.
Our experiments on seven literary and functional styles, using paired base and
aligned models (\qwenbase and \qweninst), yielded a clear answer:
alignment \emph{improves} controllable access to stylistic distributions
(style fidelity $7.26$ vs.\ $3.74$, $p < 0.0001$) rather than reducing it.
The aligned model is the only condition that reliably differentiates between
target styles (ANOVA $F = 6.32$, $p = 0.002$).

Distribution arithmetic confirms that base model distributions survive alignment
intact: logit-space interpolation produces smooth, monotonic transitions between
base and aligned behavior. However, the base model's contribution to the mixture
is not ``recovery'' of lost capability but rather dilution of the control signal.
The relevant tradeoff is not preservation versus loss, but controllability versus
distributional freedom.

These findings suggest that future work on preserving diversity during alignment
should focus less on preventing distributional loss (which does not appear to
occur for stylistic capabilities) and more on maintaining or expanding the
\emph{addressability} of latent distributions. Promising directions include
testing safety-suppressed styles where alignment may genuinely restrict access,
comparing weight-space and logit-space interpolation, measuring unsolicited
diversity (generation without style instructions), and investigating whether
lightweight steering vectors can provide the addressing mechanism without
full alignment.
